\documentclass{article}
\usepackage[T1]{fontenc}
\usepackage{lmodern}
\usepackage{graphicx}
\usepackage{amsmath,amssymb}
\usepackage[a4paper,margin=2cm]{geometry}
\usepackage[absolute,overlay]{textpos}
\setlength{\TPHorizModule}{1mm}
\setlength{\TPVertModule}{1mm}
\usepackage[indonesian]{babel}
\usepackage{hyperref}
\setlength{\emergencystretch}{2em}
% unit: Halaman 10

\begin{document}

% === Halaman 10 (llm) ===

\begin{itemize}
    \item Sifat-sifat operasi XOR:
    \begin{enumerate}[(i)]
        \item $a \oplus a = 0$
        \item $a \oplus b = b \oplus a$
        \item $a \oplus (b \oplus c) = (a \oplus b) \oplus c$
    \end{enumerate}
\end{itemize}

Contoh:
\begin{enumerate}[(i)]
    \item $1 \oplus 1 = 0$
    \item $1 \oplus 0 = 0 \oplus 1 = 1$
    \item $1 \oplus (0 \oplus 1) = (1 \oplus 0) \oplus 1 = 0$
\end{enumerate}

\end{document}
